% Gabarit pandoc des corrigés publiés (brut et guidé). Contenu sans marque ni nom de site.
% Compilation : lualatex (latexmk -lualatex).
% Le Markdown reste la source unique : ne pas éditer ce fichier .tex à la main.
\documentclass[11pt,a4paper]{article}

% --- Moteur, polices (Latin Modern) et langue ---------------------------------
\usepackage{iftex}
\ifLuaTeX\else\errmessage{Ce document se compile avec lualatex}\fi
\usepackage{amsmath,amssymb,mathtools}
\usepackage{fontspec}
% mathrm=sym : \mathrm prend ses glyphes dans la police mathématique (ℓ, μ, Ω), comme KaTeX.
\usepackage[math-style=TeX,bold-style=TeX,mathrm=sym]{unicode-math}
% Repli pour les glyphes absents de Latin Modern (indices et exposants Unicode,
% flèches, filets des schémas en texte, symboles) : DejaVu, puis Latin Modern Math
% (ℓ, lettres mathématiques Unicode des légendes).
\directlua{
  luaotfload.add_fallback("repliTexte", {"DejaVuSerif:mode=harf;", "DejaVuSans:mode=harf;", "file:latinmodern-math.otf:mode=harf;"})
  luaotfload.add_fallback("repliMono", {"DejaVuSansMono:mode=harf;", "DejaVuSans:mode=harf;", "file:latinmodern-math.otf:mode=harf;"})
}
\setmainfont{Latin Modern Roman}[RawFeature={fallback=repliTexte}]
\setsansfont{Latin Modern Sans}[RawFeature={fallback=repliTexte}]
\setmonofont{Latin Modern Mono}[RawFeature={fallback=repliMono}]
\setmathfont{Latin Modern Math}
\usepackage[french]{babel}
\frenchsetup{SmallCapsFigTabCaptions=false}

% --- Mise en page --------------------------------------------------------------
\usepackage[a4paper,margin=2.2cm,headheight=14pt,footskip=1.2cm]{geometry}
\usepackage{microtype}
\setlength{\emergencystretch}{4em}
\tolerance=1500
\setlength{\parindent}{0pt}
\setlength{\parskip}{0.55\baselineskip plus 2pt minus 1pt}
\usepackage{xcolor}
% Gris nommés : babel-french rend « ! » actif dans le document.
\definecolor{gris2}{gray}{0.98}
\definecolor{gris3}{gray}{0.97}
\definecolor{gris5}{gray}{0.95}
\definecolor{gris35}{gray}{0.65}
\definecolor{gris40}{gray}{0.60}
\definecolor{gris45}{gray}{0.55}
\definecolor{gris50}{gray}{0.50}
\definecolor{gris55}{gray}{0.45}
\definecolor{gris60}{gray}{0.40}
\definecolor{gris65}{gray}{0.35}
\definecolor{gris70}{gray}{0.30}
\definecolor{gris80}{gray}{0.20}
\definecolor{gris85}{gray}{0.15}
\usepackage{graphicx}
\usepackage{float}
\usepackage{caption}
\captionsetup{font=small,labelformat=empty,justification=centering}
\usepackage{booktabs,longtable,array,multirow,calc}
\usepackage{siunitx}
\sisetup{locale=FR}
\usepackage[version=4]{mhchem}
\usepackage{listings}
\usepackage{fancyhdr}
\usepackage{lastpage}
\usepackage[most]{tcolorbox}
\usepackage{enumitem}
\setlist{nosep,leftmargin=*,topsep=2pt}
\usepackage[hidelinks,unicode]{hyperref}
\usepackage{bookmark}
\hypersetup{pdftitle={Centrale-Supélec MP-MPI-PC-PSI — Rédaction 2026},pdfauthor={Rédigé par IA},pdfcreator={LaTeX},pdfsubject={Corrigé guidé}}
\usepackage{xurl}

% --- Pandoc : listes, tableaux, figures ----------------------------------------
\providecommand{\passthrough}[1]{#1}
\providecommand{\tightlist}{\setlength{\itemsep}{0pt}\setlength{\parskip}{0pt}}

\makeatletter
\def\maxwidth{\ifdim\Gin@nat@width>\linewidth\linewidth\else\Gin@nat@width\fi}
\def\maxheight{\ifdim\Gin@nat@height>0.55\textheight 0.55\textheight\else\Gin@nat@height\fi}
\makeatother
\setkeys{Gin}{width=\maxwidth,height=\maxheight,keepaspectratio}
\setcounter{secnumdepth}{-\maxdimen}
\providecommand{\pandocbounded}[1]{#1}
\usepackage{longtable,booktabs,array}

% --- Code (listings) -----------------------------------------------------------
\lstset{
  basicstyle=\ttfamily\small,
  columns=fullflexible,
  keepspaces=true,
  breaklines=true,
  breakatwhitespace=false,
  postbreak=\mbox{\textcolor{gris50}{$\hookrightarrow$}\space},
  frame=single,
  rulecolor=\color{gris35},
  framesep=4pt,
  xleftmargin=5pt,xrightmargin=5pt,
  backgroundcolor=\color{gris3},
  keywordstyle=\bfseries,
  commentstyle=\color{gris65},
  stringstyle=\color{gris80},
  showstringspaces=false,
  tabsize=4,
  extendedchars=true,
  aboveskip=6pt,belowskip=6pt,
}

% --- Macros des corrigés (KaTeX ↔ LaTeX) ---------------------------------------
% Les corrigés sont écrits pour KaTeX : on définit ici ce que KaTeX connaît et
% que LaTeX n'a pas, sans jamais retoucher le Markdown.
\providecommand{\R}{\mathbb{R}}
\providecommand{\N}{\mathbb{N}}
\providecommand{\Z}{\mathbb{Z}}
\providecommand{\Q}{\mathbb{Q}}
\providecommand{\C}{\mathbb{C}}
\providecommand{\dd}{\mathrm{d}}
\providecommand{\llbracket}{\lBrack}
\providecommand{\rrbracket}{\rBrack}
\providecommand{\lBrace}{\lbrace}
\providecommand{\rBrace}{\rbrace}
\providecommand{\Box}{\square}
\providecommand{\Diamond}{\lozenge}
\providecommand{\varnothing}{\emptyset}
\providecommand{\argmax}{\operatorname*{arg\,max}}
\providecommand{\argmin}{\operatorname*{arg\,min}}
\renewcommand{\leq}{\leqslant}
\renewcommand{\geq}{\geqslant}

% Équation hors texte : réduite à la largeur de la ligne si elle déborde.
\newsavebox{\boiteequation}
\newcommand{\equationajustee}[1]{%
  \sbox{\boiteequation}{$\displaystyle #1$}%
  \ifdim\wd\boiteequation>\linewidth
    \[\resizebox{\linewidth}{!}{\usebox{\boiteequation}}\]%
  \else
    \[#1\]%
  \fi}

% --- Encadrés pédagogiques (sobres, lisibles en noir et blanc) -----------------
\tcbset{
  encadre/.style={
    enhanced,breakable,sharp corners,
    colback=white,colframe=gris55,
    boxrule=0.4pt,leftrule=2.2pt,
    left=6pt,right=6pt,top=3pt,bottom=3pt,
    before skip=6pt,after skip=6pt,
    fonttitle=\bfseries\small,coltitle=black,
    attach title to upper,after title={\quad},
    parbox=false,
  }
}
\newtcolorbox{indication}[1]{encadre,colframe=gris45,colback=gris2,title={Indication~#1}}
\newtcolorbox{methode}{encadre,colframe=gris70,title={Méthode}}
\newtcolorbox{piege}{encadre,colframe=gris85,boxrule=0.8pt,leftrule=2.2pt,title={Piège}}
\newtcolorbox{jury}{encadre,colframe=gris60,colback=gris5,title={Jury}}
\newtcolorbox{encadrenote}{encadre,colframe=gris40}

% --- Titres de section ---------------------------------------------------------
\usepackage{titlesec}
\titleformat{\section}{\Large\bfseries}{}{0pt}{}
\titleformat{\subsection}{\large\bfseries}{}{0pt}{}
\titleformat{\subsubsection}{\normalsize\bfseries}{}{0pt}{}
\titlespacing*{\section}{0pt}{14pt plus 4pt minus 2pt}{6pt}
\titlespacing*{\subsection}{0pt}{10pt plus 3pt minus 2pt}{4pt}
\titlespacing*{\subsubsection}{0pt}{8pt plus 2pt minus 2pt}{2pt}

% --- En-tête et pied de page ---------------------------------------------------
\pagestyle{fancy}
\fancyhf{}
\fancyhead[L]{\small\textsc{Corrigé guidé} · Centrale-Supélec MP-MPI-PC-PSI}
\fancyhead[R]{\small Rédaction 2026}
\fancyfoot[L]{\footnotesize Corrigé rédigé par IA, non officiel}
\fancyfoot[R]{\footnotesize \thepage\,/\,\pageref*{LastPage}}
\renewcommand{\headrulewidth}{0.4pt}
\renewcommand{\footrulewidth}{0pt}
\fancypagestyle{plain}{\fancyhf{}\fancyfoot[R]{\footnotesize \thepage\,/\,\pageref*{LastPage}}\fancyfoot[L]{\footnotesize Corrigé rédigé par IA, non officiel}\renewcommand{\headrulewidth}{0pt}}

\begin{document}
\thispagestyle{plain}
\begin{center}
  {\small\scshape Corrigé guidé}\par\smallskip
  {\LARGE\bfseries Centrale-Supélec MP-MPI-PC-PSI — Rédaction 2026\par}\medskip
  {\large Centrale-Supélec · MP-MPI-PC-PSI · Rédaction · 2026\par}
\end{center}
\vspace{2pt}\hrule\vspace{8pt}

\emph{Corrigé rédigé par IA, non officiel. Signalez toute erreur.}

\hypertarget{pruxe9sentation-du-sujet}{%
\section{Présentation du sujet}\label{pruxe9sentation-du-sujet}}

\textbf{Format.} Épreuve de 4 heures, sans document ni appareil électronique, en deux exercices notés ensemble (« Résumé et dissertation ont la même notation et forment un ensemble indissociable ») : on ne sacrifie ni l'un ni l'autre. Le résumé se présente en premier, la dissertation en second, en écrivant une ligne sur deux ; la présentation, l'orthographe et la clarté comptent dans la note. Un partage raisonnable : 1 h 15 pour le résumé, 2 h 30 pour la dissertation, 15 minutes de relecture.

\textbf{Partie A, résumé.} Un extrait d'environ 1 200 mots de François Jacob, \emph{Le Jeu des possibles} (1987), à réduire à 200 mots (tolérance de 10 \%, donc 180 à 220), avec une barre nette toutes les cinquantaines de mots et le total exact à la fin. Le texte est argumentatif et procède par retournements. Le mythe semble l'emporter sur la science pour l'unité et la cohérence ; mais la science, en se limitant, trouve des réponses générales, au prix d'un savoir en îlots. Les deux reposent sur le même principe, expliquer le visible par l'invisible, et commencent tous deux par imaginer un monde possible. Le mythe s'en tient là, la science soumet son rêve à l'expérience. Il reste au mythe une fonction morale que la science ne remplit pas.

\textbf{Partie B, dissertation.} Le sujet reprend une phrase du texte : « Mythique ou scientifique, la représentation du monde que construit l'homme fait toujours une large part à son imagination. » Il faut la discuter à la lumière du thème « Expériences de la nature » et des trois œuvres : \emph{Vingt mille lieues sous les mers} (Verne), \emph{Le Mur invisible} (Haushofer), \emph{La Connaissance de la vie} (Canguilhem). Contraintes : confronter les trois œuvres et y renvoyer avec précision, deux ou trois parties, 1 800 mots au plus, tout dépassement manifeste étant sanctionné.

\textbf{Ce qui fait la difficulté.} Le texte du résumé et le sujet de dissertation se recoupent : il faut résumer Jacob en A, puis s'en servir sans le résumer une seconde fois en B. La formule ressemble à une évidence (« on imagine partout ») ; tout son intérêt se loge dans deux mots, « toujours » et « une large part ».

\textbf{À ne pas laisser passer.} Les barres et le total exact du résumé (vérifiables en quelques secondes), l'annonce du plan, la présence des trois œuvres dans chaque partie, le respect du plafond de mots.

\hypertarget{partie-a-ruxe9sumuxe9-de-texte}{%
\section{Partie A --- Résumé de texte}\label{partie-a-ruxe9sumuxe9-de-texte}}

Il s'agit de réduire le texte à un sixième de sa longueur, en gardant son ordre, son énonciation (on écrit à la place de l'auteur, sans « l'auteur dit ») et ses articulations logiques.

\hypertarget{uxe9tape-1-lire-le-texte-et-duxe9gager-la-thuxe8se}{%
\subsection{Étape 1 --- Lire le texte et dégager la thèse}\label{uxe9tape-1-lire-le-texte-et-duxe9gager-la-thuxe8se}}

\begin{indication}{1}

Repérez les connecteurs qui retournent la perspective (« À première vue », « Cependant », « Mais », « au contraire », « En outre ») : ils découpent le texte.

\end{indication}

\begin{indication}{2}

Demandez-vous ce que mythe et science ont en commun selon Jacob, puis ce qui les sépare : la thèse tient dans cette double réponse.

\end{indication}

\hypertarget{solution}{%
\subsubsection{Solution}\label{solution}}

Le texte ne place pas la science au-dessus du mythe, ni l'inverse. Il les rapproche puis les distingue. Ce qu'ils ont en commun : expliquer le monde visible par des forces invisibles, et donc commencer par \textbf{l'invention d'un monde possible}, ce qui donne à l'imagination une large part dans toute représentation du monde. Ce qui les sépare : le mythe s'arrête à cette invention et y fait entrer la réalité, la science \textbf{expose sans relâche son rêve à la critique et à l'expérience}. Deux paradoxes encadrent la thèse. Au début, le mythe l'emporte en unité, mais la science gagne la généralité en posant des questions limitées. À la fin, le mythe garde une fonction propre, donner sens, valeurs et loi, que la science, qui veut une connaissance sans émotion, ne remplit pas.

\begin{piege}

Lire le texte comme un éloge de la science contre la superstition. Jacob ne dit nulle part que la science « est vraie » et le mythe « faux » : il oppose deux façons de traiter l'imagination, et il termine en reconnaissant au mythe ce que la science ne donne pas.

\end{piege}

\hypertarget{uxe9tape-2-duxe9couper-le-texte-et-ruxe9partir-les-mots}{%
\subsection{Étape 2 --- Découper le texte et répartir les mots}\label{uxe9tape-2-duxe9couper-le-texte-et-ruxe9partir-les-mots}}

\begin{indication}{1}

Regroupez les sept paragraphes en quelques mouvements de sens, puis donnez à chacun un nombre de mots à peu près proportionnel à sa longueur.

\end{indication}

\begin{indication}{2}

Les paragraphes 4 et 5 forment un seul mouvement : le principe commun, puis sa conséquence sur ce qu'est observer.

\end{indication}

\hypertarget{solution-1}{%
\subsubsection{Solution}\label{solution-1}}

\begin{longtable}[]{@{}
  >{\raggedright\arraybackslash}p{(\columnwidth - 6\tabcolsep) * \real{0.1205}}
  >{\raggedright\arraybackslash}p{(\columnwidth - 6\tabcolsep) * \real{0.1333}}
  >{\raggedright\arraybackslash}p{(\columnwidth - 6\tabcolsep) * \real{0.5891}}
  >{\raggedright\arraybackslash}p{(\columnwidth - 6\tabcolsep) * \real{0.1271}}@{}}
\toprule\noalign{}
\begin{minipage}[b]{\linewidth}\raggedright
Mouvement
\end{minipage} & \begin{minipage}[b]{\linewidth}\raggedright
Paragraphes
\end{minipage} & \begin{minipage}[b]{\linewidth}\raggedright
Idée
\end{minipage} & \begin{minipage}[b]{\linewidth}\raggedright
Mots visés
\end{minipage} \\
\midrule\noalign{}
\endhead
\bottomrule\noalign{}
\endlastfoot
1 & §1 à §3 & Le mythe l'emporte en unité ; la science, locale, est née des questions limitées, qui mènent aux réponses générales ; savoir choisir ses problèmes fait le chercheur créatif ; savoir en îlots, unifications partielles & ≈ 85 \\
2 & §4 et §5 & Même principe, le visible expliqué par l'invisible ; observer suppose une idée préalable ; toute enquête commence par l'invention d'un monde possible & ≈ 55 \\
3 & §6 & Le mythe s'en tient là et lit chaque fait comme un signe ; la science confronte sans relâche son rêve au réel & ≈ 35 \\
4 & §7 & Fonctions du mythe (angoisse, sens, valeurs, loi) ; la science veut une connaissance sans émotion & ≈ 25 \\
\end{longtable}

Les exemples sont sacrifiés : les questions sur la pierre et le sang, les trois unifications historiques (mécanique statistique, Maxwell, mécanique quantique), la foudre de Zeus, le microbe, le nom de Medawar.

\begin{methode}

Un budget de mots par mouvement, fixé avant d'écrire, évite le défaut classique du résumé : développer le début et bâcler la fin.

\end{methode}

\hypertarget{uxe9tape-3-reformuler-sans-coller-ni-ajouter}{%
\subsection{Étape 3 --- Reformuler sans coller ni ajouter}\label{uxe9tape-3-reformuler-sans-coller-ni-ajouter}}

\begin{indication}{1}

Pour chaque mouvement, écrivez une phrase-noyau avec vos propres mots ; ne gardez une expression de Jacob que si c'est un concept que rien ne remplace.

\end{indication}

\begin{indication}{2}

Le paradoxe des questions limitées et le partage clôture/confrontation sont les deux idées les plus denses : réservez-leur une formulation nette.

\end{indication}

\hypertarget{solution-2}{%
\subsubsection{Solution}\label{solution-2}}

On reformule « unité et cohérence » (besoin de voir le monde unifié), « émiettement » et « îlots séparés » (le savoir morcelé), « s'expose à la critique et à l'expérience » (soumettre son rêve à la critique et au contrôle expérimental). « Monde possible » est gardé : c'est le concept du texte. On s'interdit tout jugement absent du texte : la dernière phrase de Jacob constate que la science veut se libérer de l'émotion ; elle ne l'en loue ni ne le lui reproche.

\begin{piege}

Résumer le paragraphe 6 par « la science vérifie » : on perd l'idée que la science \textbf{aussi} rêve. C'est ce qui la rapproche du mythe et donne son sens au partage.

\end{piege}

\hypertarget{uxe9tape-4-ruxe9diger-compter-placer-les-barres}{%
\subsection{Étape 4 --- Rédiger, compter, placer les barres}\label{uxe9tape-4-ruxe9diger-compter-placer-les-barres}}

\begin{indication}{1}

Rédigez d'abord au brouillon, comptez, puis coupez les mots-outils et les redites plutôt que les idées.

\end{indication}

\begin{indication}{2}

Choisissez une convention de comptage et appliquez-la partout. Le plus sûr est la plus stricte, où chaque élision compte pour un mot ; un texte sans élision ni trait d'union donne le même total quelle que soit la convention.

\end{indication}

\hypertarget{solution-3}{%
\subsubsection{Solution}\label{solution-3}}

Le mythe satisfait mieux que la science notre besoin de voir le monde unifié : il explique tout a priori ; la science, locale, se contente de réponses partielles et provisoires. Mais elle est née en substituant aux questions générales des questions limitées, qui, paradoxalement, mènent aux réponses générales. Choisir les problèmes \textbar{} mûrs fait le chercheur créatif ; le débutant veut tout embrasser. Cette méthode morcelle le savoir en îlots, que des synthèses relient parfois sans combler toutes les brèches.

Mythe et science obéissent pourtant au même principe : expliquer le visible par des forces invisibles. Tous deux imaginent. Observer suppose de savoir déjà \textbar{} quoi examiner ; les découvertes naissent souvent moins des instruments que de regards neufs, guidés par une idée de la réalité. Toute enquête commence par inventer un monde possible.

Le mythe en reste là : il loge le réel dans son cadre et lit chaque fait comme un signe. La science, elle, \textbar{} soumet sans relâche son rêve à la critique et au contrôle expérimental, pour approcher davantage la réalité.

Reste au mythe une fonction que la science, soucieuse de soustraire la connaissance à toute émotion, ne remplit pas : apaiser notre angoisse, donner sens à la vie, fonder des valeurs et une loi. \textbar{}

\textbf{Total : 200 mots.}

Les barres tombent après le cinquantième mot (« problèmes »), le centième (« déjà »), le cent-cinquantième (« elle, ») et le deux-centième. Le résumé ne contient ni élision ni trait d'union : le total ne dépend donc pas de la convention retenue. « A priori » compte pour deux mots.

\begin{piege}

Annoncer un total arrondi, ou dérivé d'une convention qu'on n'applique pas jusqu'au bout. La tolérance de 10 \% ne dispense pas d'annoncer le total \textbf{exact}, et une barre mal placée se voit immédiatement.

\end{piege}

\begin{jury}

Dans une épreuve de résumé, on attend une reformulation personnelle qui garde l'ordre et l'énonciation du texte, sans exemple ni commentaire, et un décompte honnête. Un décompte inexact est sanctionné même quand le nombre de mots reste dans la tolérance.

\end{jury}

\hypertarget{partie-b-dissertation}{%
\section{Partie B --- Dissertation}\label{partie-b-dissertation}}

Il faut discuter une phrase que l'on vient de résumer, mais en l'éprouvant sur les trois œuvres, sans paraphraser Jacob.

\hypertarget{uxe9tape-1-analyser-la-formule}{%
\subsection{Étape 1 --- Analyser la formule}\label{uxe9tape-1-analyser-la-formule}}

\begin{indication}{1}

Pesez chaque terme : « représentation », « construit », « une large part », « toujours », « mythique ou scientifique ».

\end{indication}

\begin{indication}{2}

Repérez la restriction et l'universalité que la phrase associe : la discussion naîtra de leur tension.

\end{indication}

\hypertarget{solution-4}{%
\subsubsection{Solution}\label{solution-4}}

« Représentation » et « construit » : notre rapport à la nature n'est pas un enregistrement, mais une construction. « Mythique ou scientifique » : la phrase range sous une même loi deux régimes qu'on oppose d'ordinaire. « Toujours » : la thèse est universelle et vaut donc pour la science, qu'on croit volontiers tirée de la seule observation. « Une large part » : la restriction est essentielle ; l'imagination n'est pas tout, il y a un reste (l'expérience, la résistance des faits). Dans le texte de Jacob, la phrase sert à rapprocher le mythe et la science avant de les distinguer par ce qu'ils font de l'imaginaire.

\begin{piege}

Deux contresens opposés : réfuter la formule (« la science, elle, ne rêve pas »), ou l'illustrer platement (« on imagine partout »). Dans les deux cas on a oublié « une large part ».

\end{piege}

\hypertarget{uxe9tape-2-construire-la-probluxe9matique}{%
\subsection{Étape 2 --- Construire la problématique}\label{uxe9tape-2-construire-la-probluxe9matique}}

\begin{indication}{1}

Si l'imaginaire est partout, que devient la différence entre connaître et fabuler ?

\end{indication}

\begin{indication}{2}

Le thème porte sur des \emph{expériences} de la nature : demandez-vous si toute expérience se ramène à une représentation.

\end{indication}

\hypertarget{solution-5}{%
\subsubsection{Solution}\label{solution-5}}

Deux questions liées : si le savant imagine comme le conteur, qu'est-ce qui distingue encore la connaissance de la fable ? Et l'expérience de la nature se réduit-elle à la représentation que nous nous en faisons ? La première question permet de nuancer la formule, la seconde de la dépasser : elle oblige à chercher dans les œuvres ce que la formule laisse de côté.

\hypertarget{uxe9tape-3-buxe2tir-le-plan-et-ruxe9partir-les-ruxe9fuxe9rences}{%
\subsection{Étape 3 --- Bâtir le plan et répartir les références}\label{uxe9tape-3-buxe2tir-le-plan-et-ruxe9partir-les-ruxe9fuxe9rences}}

\begin{indication}{1}

Trois temps : confirmer la formule, la nuancer à partir de sa propre restriction, puis la dépasser à partir de ce qu'elle ne dit pas.

\end{indication}

\begin{indication}{2}

Pour chaque partie, cherchez dans chaque œuvre une scène ou une étude précise. Canguilhem se cite par le titre de ses études, pas en le « racontant ».

\end{indication}

\hypertarget{solution-6}{%
\subsubsection{Solution}\label{solution-6}}

\begin{longtable}[]{@{}
  >{\raggedright\arraybackslash}p{(\columnwidth - 6\tabcolsep) * \real{0.2053}}
  >{\raggedright\arraybackslash}p{(\columnwidth - 6\tabcolsep) * \real{0.2288}}
  >{\raggedright\arraybackslash}p{(\columnwidth - 6\tabcolsep) * \real{0.3151}}
  >{\raggedright\arraybackslash}p{(\columnwidth - 6\tabcolsep) * \real{0.2208}}@{}}
\toprule\noalign{}
\begin{minipage}[b]{\linewidth}\raggedright
Partie
\end{minipage} & \begin{minipage}[b]{\linewidth}\raggedright
Verne
\end{minipage} & \begin{minipage}[b]{\linewidth}\raggedright
Canguilhem
\end{minipage} & \begin{minipage}[b]{\linewidth}\raggedright
Haushofer
\end{minipage} \\
\midrule\noalign{}
\endhead
\bottomrule\noalign{}
\endlastfoot
I. L'imagination au principe de toute représentation & « Le pour et le contre » : Aronnax infère un narval gigantesque qu'il n'a jamais vu & « Machine et organisme », « La théorie cellulaire », « Le vivant et son milieu », « La monstruosité et le monstrueux » : origine imaginaire et technique des concepts & Hypothèses sur le mur ; le choix d'écrire un rapport \\
II. Ce qui sépare science et mythe : l'épreuve infligée à l'imaginaire & Tôles boulonnées, mesures, le pôle Sud, « Les poulpes » ; contre-exemple de l'Atlantide & « L'expérimentation en biologie animale » : le vivant résiste au protocole sans l'abolir & Fenaison, pommes de terre, naissance du veau : des images défaites par le travail \\
III. Une expérience qui déborde la représentation & Nemo, savant et homme d'un mythe personnel (« la mer est tout ») & « La pensée et le vivant », le vitalisme comme exigence, le milieu propre du savant & Le mur inexpliqué, la mort de Lynx et du taureau, la corneille blanche \\
\end{longtable}

\begin{methode}

Faites un tableau croisé parties × œuvres avant de rédiger : une case vide signale une partie qui ne confronte pas les trois œuvres.

\end{methode}

\begin{jury}

Un plan qui consacre une partie à chaque œuvre ne répond pas à la consigne de confrontation, quelle que soit sa qualité. On attend aussi des renvois situés (chapitre, scène, titre d'étude) plutôt que des allusions générales.

\end{jury}

\hypertarget{uxe9tape-4-ruxe9diger-lintroduction}{%
\subsection{Étape 4 --- Rédiger l'introduction}\label{uxe9tape-4-ruxe9diger-lintroduction}}

\begin{indication}{1}

Partez d'une scène concrète d'une œuvre qui montre déjà la tension entre image et réalité, puis citez la formule en entier.

\end{indication}

\begin{indication}{2}

Enchaînez dans cet ordre : analyse des termes, problématique, présentation des œuvres, annonce explicite du plan.

\end{indication}

\hypertarget{solution-7}{%
\subsubsection{Solution}\label{solution-7}}

On ouvre sur la chasse au « monstre » de l'\emph{Abraham Lincoln} : une frégate entière poursuit une créature composée par la presse. Puis vient la citation, l'analyse des trois affirmations (construction, « une large part », « toujours »), les deux questions, la présentation en une phrase de chaque œuvre, et l'annonce des trois parties. Le texte complet est donné à l'étape 6.

\begin{piege}

Oublier d'annoncer le plan. Une introduction qui pose bien la question mais ne dit pas le chemin oblige le correcteur à le deviner.

\end{piege}

\hypertarget{uxe9tape-5-conclure}{%
\subsection{Étape 5 --- Conclure}\label{uxe9tape-5-conclure}}

\begin{indication}{1}

Répondez nettement à la problématique : dans quelle mesure la formule est-elle juste ?

\end{indication}

\begin{indication}{2}

Reprenez une référence par œuvre en une ligne, puis terminez sur ce que la troisième partie a ajouté, sans ouvrir une question artificielle.

\end{indication}

\hypertarget{solution-8}{%
\subsubsection{Solution}\label{solution-8}}

La formule est juste, mais incomplète. Elle est juste parce qu'aucune représentation de la nature ne se passe d'invention (le narval d'Aronnax, la généalogie des concepts chez Canguilhem, les hypothèses de la narratrice). La science se distingue du mythe par l'épreuve qu'elle impose à cette invention, que le vivant complique sans l'abolir. Elle est incomplète, parce qu'elle ne dit rien d'une expérience de la nature qui précède toute représentation, celle que donne Haushofer : la nature n'est pas seulement ce que nous imaginons, mais ce qui nous arrive.

\hypertarget{uxe9tape-6-le-devoir-complet}{%
\subsection{Étape 6 --- Le devoir complet}\label{uxe9tape-6-le-devoir-complet}}

\begin{indication}{1}

Chaque paragraphe du développement développe une seule idée, s'appuie sur une œuvre et s'achève sur ce qu'il apporte à la discussion ; une transition sous forme de question relie les parties.

\end{indication}

\begin{indication}{2}

Visez 1 400 à 1 600 mots et comptez au fil de la rédaction, par paragraphe : le plafond de 1 800 ne se découvre pas à la relecture.

\end{indication}

\hypertarget{solution-9}{%
\subsubsection{Solution}\label{solution-9}}

\textbf{Problématique :} si le savant imagine comme le conteur, qu'est-ce qui distingue encore la connaissance de la fable, et l'expérience de la nature se réduit-elle à la représentation que nous nous en faisons ?

Au début de \emph{Vingt mille lieues sous les mers}, une frégate entière, l'\emph{Abraham Lincoln}, appareille pour chasser un être qui n'existe pas : le « monstre » que la presse a composé à partir de quelques abordages inexpliqués. Avant d'être observée, la mer est déjà peuplée d'images. C'est ce que donne à penser François Jacob dans \emph{Le Jeu des possibles} : « Mythique ou scientifique, la représentation du monde que construit l'homme fait toujours une large part à son imagination. » La phrase réunit trois affirmations. Le monde ne nous est pas simplement donné : nous en \emph{construisons} une représentation. L'imagination y occupe « une large part », ce qui veut dire qu'elle n'occupe pas toute la place. Enfin, cela vaut « toujours », donc aussi pour la science, que l'on croit d'ordinaire tirée de la seule observation et opposée trait pour trait au mythe. Si le savant imagine comme le conteur, qu'est-ce qui distingue encore la connaissance de la fable ? Et l'expérience de la nature se réduit-elle à la représentation que nous nous en faisons ? Le roman d'aventures scientifiques de Jules Verne, le récit de survie de Marlen Haushofer et les études de Georges Canguilhem sur la biologie permettent de mettre la formule à l'épreuve. Nous montrerons d'abord que l'imagination est au principe de toute représentation de la nature, savante comprise ; puis que ce qui sépare la science du mythe tient au traitement que l'on fait subir à cette part imaginaire ; enfin, que l'expérience de la nature déborde la représentation elle-même et rend à l'imaginaire une autre fonction que celle de connaître.

Nulle représentation de la nature, d'abord, ne commence par la seule observation. Le roman de Verne s'ouvre sur une énigme et sur un concours d'hypothèses. Dans le chapitre « Le pour et le contre », Aronnax, professeur au Muséum, écarte l'écueil flottant et l'engin sous-marin, puis conclut à un narval gigantesque armé d'une défense capable de percer une coque. Il n'a rien vu : il raisonne sur la déchirure de la \emph{Scotia} et sur des témoignages. Son savoir lui sert à donner une forme plausible à ce qu'il imagine, et Verne souligne l'ironie : l'hypothèse écartée, celle d'une machine construite par un homme, est la vraie. Le savant n'a pas été moins imaginatif que la foule ; il l'a été avec plus de méthode.

Canguilhem montre que cette part imaginaire ne se limite pas aux rumeurs : elle loge au cœur des concepts. Dans « Machine et organisme », il établit que l'explication mécaniste du vivant ne sort pas de l'observation des corps, mais des techniques disponibles, automates et horloges, dont on transporte le modèle dans l'organisme. « La théorie cellulaire » rappelle que la cellule de Hooke est d'abord l'alvéole d'un rayon de miel, et que la notion a longtemps porté des valeurs affectives et sociales de coopération et d'association, qui ont orienté la recherche autant que les microscopes. « Le vivant et son milieu » retrace l'importation en biologie d'une notion venue de la mécanique. Enfin, « La monstruosité et le monstrueux » rappelle que le monstre fut une créature de la légende avant de devenir l'objet d'une science, la tératologie : l'imaginaire précède le concept et lui lègue ses questions.

La narratrice du \emph{Mur invisible} n'échappe pas à cette loi. Devant la paroi transparente qui la sépare du reste du monde, elle échafaude des explications, l'effet de quelque arme inconnue par exemple, sans pouvoir en vérifier aucune. Puis elle entreprend d'écrire un rapport : cette forme sobre, chronologique, est déjà une manière de construire un monde habitable, pour ne pas perdre la raison. L'imagination n'est donc pas un ornement de la connaissance : elle en est la condition. Mais Jacob ne dit pas qu'elle est tout, et c'est sur ce reste qu'il faut s'arrêter.

« Une large part » suppose en effet une autre part, celle du réel qui répond. Le texte de Jacob situe la différence non dans la présence de l'imagination, mais dans ce qu'on en fait : le mythe s'arrête au monde qu'il a inventé, la science confronte sans relâche le possible et l'existant. Verne met en scène cette confrontation. Le monstre devient une coque de tôles boulonnées ; à bord du \emph{Nautilus}, Aronnax mesure, sonde, relève la position, compte les espèces derrière les vitres du salon. Au pôle Sud, Nemo attend que le soleil paraisse pour prendre sa hauteur et prouver, par le calcul, qu'il a atteint le point visé. Le chapitre « Les poulpes » accomplit tout le trajet : la légende du kraken, évoquée par les voyageurs, n'est ni confirmée ni abolie ; un calmar géant surgit, et l'on peut décrire ses bras, son bec, sa taille. La fable avait désigné, en l'exagérant, un possible que l'expérience corrige. Le roman sait pourtant aussi franchir la frontière dans l'autre sens : au chapitre « Un continent disparu », Nemo trace un seul mot sur un rocher, « Atlantide », et la légende de Platon devient paysage sans autre examen. Le romancier s'accorde la liberté que le savant se refuse, et c'est sa propre fiction qui occupe alors toute la place.

Canguilhem empêche de croire que cette épreuve soit un tribunal simple. « L'expérimentation en biologie animale » énumère ce qui, dans le vivant, résiste aux protocoles calqués sur la physique : la spécificité, qui interdit de conclure d'une espèce à une autre, l'individualité, la totalité de l'organisme, que l'on altère en l'isolant, et l'irréversibilité des phénomènes. Le biologiste doit donc inventer des méthodes ajustées à son objet ; l'imagination n'est pas éliminée, elle est corrigée et relancée. Mais Canguilhem ne conclut jamais que les concepts équivalent à des mythes : la cellule, chargée d'images au départ, est devenue un concept opératoire parce qu'elle a été mise à l'épreuve.

Haushofer donne de cette correction la version la plus quotidienne. La citadine arrive dans la montagne avec des images de vacances ; l'hiver, la fenaison dans la clairière, la culture des pommes de terre, la naissance du veau de Bella lui imposent un savoir qui ne se discute pas. Elle apprend en se trompant, elle se méfie de ses rêveries et de ses souvenirs, car il lui faut d'abord survivre. La forêt n'est ni bienveillante ni hostile : elle obéit à son propre rythme. Ici la correction ne réfute pas des théories, elle défait des images. Faut-il pour autant tenir l'imagination pour une illusion à réduire, que la science effacerait peu à peu ?

Ce serait oublier la fonction que Jacob reconnaît lui-même au mythe : donner un sens, des valeurs, une loi, aider à supporter l'angoisse. Canguilhem prolonge cette idée en biologiste. Dès « La pensée et le vivant », la connaissance apparaît comme une manière, pour le vivant humain, de résoudre les tensions qui l'opposent à son milieu ; elle n'est pas une pure contemplation. S'il défend le vitalisme, c'est comme une exigence plutôt que comme une doctrine scientifique : la fidélité de la pensée à l'originalité de la vie. Et « Le vivant et son milieu » rappelle que l'univers abstrait du savant reste ancré dans le milieu propre d'un vivant. La part imaginaire de la représentation exprime alors un rapport vital au monde, non une faiblesse de l'esprit.

Verne l'incarne en Nemo. Avec sa bibliothèque de douze mille volumes et ses collections, nul n'est plus savant ; pourtant « la mer est tout » est une profession de foi, non une loi physique. Sa mer est un refuge et une patrie, le lieu d'une liberté et d'une vengeance. Aronnax, devant les panneaux ouverts du salon, oublie qu'il est prisonnier, et les énumérations de genres et d'espèces deviennent chez Verne une poésie de la nomenclature. Le même livre fait coexister les deux parts sans les confondre.

Le \emph{Mur invisible} va plus loin, et touche la limite de la formule. Le mur n'est jamais expliqué : il n'est ni mythe ni objet de science, il est ce qui résiste à toute représentation. Autour de lui, la nature cesse d'être une image pour devenir le froid, la faim, la fatigue, puis la mort de Lynx et du jeune taureau, tués à la hache par un homme surgi de nulle part. Ces épreuves ne font pas une large part à l'imagination : elles la réduisent au silence. Et pourtant le récit s'achève sur un geste, aller nourrir la corneille blanche que les autres rejettent, qui ne doit rien au mythe ni à la théorie et fonde une valeur à même l'expérience. Le rapport qu'elle écrit n'est ni fable ni traité : c'est un témoignage, dont la vérité tient à la fidélité.

La formule de Jacob est donc juste, mais incomplète. Juste, parce qu'aucune représentation de la nature ne se passe d'invention : le narval d'Aronnax, les images dont Canguilhem fait la généalogie, les hypothèses de la narratrice le confirment. Mais la science se distingue du mythe par l'épreuve qu'elle impose à cette invention, une épreuve que le vivant complique sans l'abolir. Surtout, la formule ne dit rien d'une expérience de la nature qui précède toute représentation et à laquelle toute représentation doit répondre : celle qu'Haushofer donne à lire, et qui rappelle que la nature n'est pas seulement ce que nous imaginons, mais ce qui nous arrive.

\textbf{Environ 1 470 mots} (1 588 si chaque élision compte pour un mot) ; plafond : 1 800 mots.

\begin{methode}

La troisième partie ne répète pas la deuxième : elle cherche dans une œuvre ce que la formule laisse de côté. C'est là que la dissertation discute vraiment le sujet au lieu de reformuler le texte source.

\end{methode}

\begin{piege}

Attribuer à Canguilhem une thèse relativiste (« la science n'est qu'un récit ») ou faire du mur de Haushofer une allégorie plaquée. Canguilhem montre l'origine imaginaire des concepts sans conclure qu'ils valent des mythes ; le roman laisse le mur inexpliqué, et c'est ce refus qui compte.

\end{piege}

\hypertarget{bilan}{%
\section{Bilan}\label{bilan}}

\begin{itemize}
\tightlist
\item
  \textbf{Résumé} : découper le texte en mouvements et répartir les mots avant d'écrire ; garder les articulations logiques et sacrifier les exemples ; reformuler sans ajouter de jugement ; compter avec une convention stricte, placer les barres et annoncer le total exact.
\item
  \textbf{Dissertation} : analyser chaque terme de la formule, en particulier les restrictions (« une large part ») et les universels (« toujours ») ; poser une problématique qui permette à la fois de nuancer et de dépasser ; confronter les trois œuvres dans chaque partie, avec des renvois situés ; annoncer le plan ; tenir le plafond de mots.
\item
  \textbf{À revoir sur les œuvres} : chez Verne, les chapitres d'hypothèses et de vérification (« Le pour et le contre », « Le pôle Sud », « Les poulpes ») et le personnage de Nemo ; chez Canguilhem, la structure du recueil (méthode, histoire, philosophie) et la thèse de chaque étude ; chez Haushofer, le travail des saisons, les animaux et la fin du récit.
\item
  \textbf{Transposable} : le couple imagination/expérience revient dans tout sujet sur la connaissance de la nature (observation, hypothèse, modèle, rapport vécu à la nature). Le schéma « confirmer, nuancer par la restriction du sujet, dépasser par ce qu'il omet » s'applique à toute citation qui associe un universel et une restriction.
\end{itemize}

\end{document}
